\section{Homotopy Type Theory}

Homotopy Type Theory (HoTT) extends MLTT by introducing a topological/homotopical interpretation
of types, an interpretation of 

\subsection{A Homotopical Intepretation}
In HoTT, types can be interpretated as either topologically as spaces, 
or as $\infty$-groupoids from higher-dimensional category theory.
Any topological space can be classified upto homotopy equivalence by considering its \textit{fundamental
$\infty$-groupoid}. This follows from the fundamental $\infty$-groupoid construction being adjoint 
to the geometric realization functor. (This is the \textit{homotopy hypothesis/theorem}).[Pursuing Stacks]

\subsection{Identity Types}
The precise manner in which types correspond to spaces or $\infty$-groupoids is demonstrated by
constructing an associated $\infty$-groupoid.

In the homotopy interpretation, the identity type $\mathsf{Id}_A(x,y)$ corresponds to the space of 
paths from $a : A$ to $b : B$ endowed with the compact-open topology. $\refl_x$ is the
constant path at $x$.

\begin{proposition}
  Given $x, y, z : A$, there exist the following maps 
  \begin{itemize}
    \item [(i)] $(x = y) \to (y = x)$
    \item [(ii)] $(x = y) \to (y = z) \to (x = z)$
  \end{itemize}
\end{proposition}
\textit{Proof.} (i) By path induction, assume $x \deq y$ and $p \deq \refl_x$. Now simply exhibit $\mathsf{id}_{x =_A x} \deq (\lambda a.a)$. 
(ii) By path induction, assume $x \deq y$ and $p \deq \refl_x$. Exhibit $(\lambda a.b.b)$. \qed

Denote the map in $(i)$ as $(-)^{-1}$, and $(ii)$ as $(- \sqcdot -)$. Then we have the following results.

\begin{proposition} Let $x, y, z, w : A$, $p : x =_A y, q : y =_A z$ and $r : z =_A w$. Then we have:
  \begin{itemize}
    \item [(i)] $p = \refl_x \sqcdot p$, and $p = p \sqcdot \refl_y$
    \item [(ii)] $p \sqcdot p^{-1} = \refl_x$ and $p^{-1} \sqcdot p = \refl_y$
    \item [(iii)] $(p \sqcdot q) \sqcdot r = p \sqcdot (q \sqcdot r)$ 
   \end{itemize} 
\end{proposition}
\textit{Proof.} Follows from path induction. \qed

A remarkable feature is that while (i) and (ii) can be proven judgementally, (iii) is only true propositionally,
i.e., upto a path in the higher path space.
Strikingly, we also have all the required higher coherences propositionally using only path induction.
This gives us the (weak) $\infty$-groupoid interpretation. 

\subsection{Axioms of Function Extensionality and Univalence}
Some expected semantic behaviour does not directly follow from the machinery of HoTT laid out so far. To enforce such behaviour
we introduce a few `axioms' into our theory, which are merely inhabitants of some type without an inductive construction.

The notion of equality for functions is too restrictive to capture their expected behaviour.
Classically, we expect two functions to be equal if they take the same values on all the inputs.
To ensure semantic well-behavedness, we must enforce this as an axiom.

\begin{axiom}[Function Extensionality]
   Let $f, g : \Prod{x:A}{B(a)}$, then we have
   \[     (f = g) \eqv \left( \Prod{x:A}{f(x) =_{B(x)} g(x)} \right)    \]
\end{axiom}
In, Voevodsky proposed the univalence axiom, \textit{}
\begin{axiom}[Univalence]
  For $A, B : \U$, there is an equivalence:
  \[(A =_{\U} B) \eqv (A \eqv B)\]
\end{axiom}
In particular, if $A \eqv B$, then $A =_{\U} B$.
Roughly speaking, this means ``isomorphic ojects are equal''. 
In a univalent universe $\U$, isomrphic objects/equivalent types can be treated as equal and one can transport 
(theorems/proofs/properties) across them.

\subsection{Higher Inductive Types}

Following this interpretation, HoTT allows us to extend 
our domain of inductive definitions, from standard $W$-types which are types generated freely by 
its constructors, to types generated by constructors along with relations.

We allow the constructors to also be higher-dimensional paths between points. In this manner, 
free generation is equivalent to generation on points quotiented by relations (corresponding to these higher paths).

Such a \textit{higher inductive type} hardcoded in HoTT is the circle type $\Sph^1$ generated by the constructors
\begin{itemize}
  \item A point $\base : \Sph^1$
  \item A path $ \mathsf{loop} : \base =_{\Sph^1} \base $
\end{itemize}
This inductive type corresponds to a chosen basepoint $\base$, and a nontrivial path $\mathsf{loop}$ connected the basepoint,
giving topologically a circle.

One can build many of the familiar mathematical objects such as the integers, higher spheres, CW-complexes etc. using
this general schema. Various approaches exist to formalize such a schema, but have only yeilded partial satisfaction.
Some of the existing approaches are using Quotient Inductive Types [], Homotopy Initial Algebras [], and computationally
most satisfactory, Cubical Type Theory []. A satisfactory formalism remains an open problem.
